\documentclass[10pt,a4paper]{amsart}
\usepackage[margin=19mm]{geometry}
\usepackage{amsmath,amssymb,mathtools}
\usepackage[scr=rsfs,cal=euler]{mathalfa}
\usepackage{xfrac}
\usepackage[unicode,hidelinks,bookmarks=false]{hyperref}
\newcommand{\ie}{i.e.\ }
\newcommand{\cf}{cf.\ }
\newcommand{\resp}{respectively\ }
\newcommand{\Subsection}{\subsection}
\providecommand{\nobreakdash}{\nobreak\hbox{-}}
\setlength{\emergencystretch}{2em}
\multlinegap0pt
\usepackage{amssymb}
\usepackage{enumerate}
\usepackage{xcolor}
\usepackage{tikz}
\usepackage[shortlabels]{enumitem}
\newtheorem{prop}{Proposition}
\newtheorem{lemma}{Lemma}
\newcommand{\A}{\mathcal{A}}
\newcommand{\p}{\partial}
\title{The Hard Lefschetz Theorem on K\"ahler Lie Algebroids}
\author{Shane Rankin}
\date{Source: arXiv:2504.12550v2 (CC BY 4.0)}
\begin{document}
\maketitle

\noindent\emph{Source and licence.} The Hard Lefschetz Theorem on K\"ahler Lie Algebroids. Version arXiv:2504.12550v2 (CC BY 4.0), distributed by arXiv under the Creative Commons Attribution 4.0 International licence, http://creativecommons.org/licenses/by/4.0/, as recorded in the arXiv licence metadata for this exact version and shown on the abstract page https://arxiv.org/abs/2504.12550v2. Full licence notice: \texttt{licenses/arxiv-2504.12550-CC-BY-4.0.txt}.\par\smallskip

\noindent\textit{Editorial scope.} Counted unit: the article's lemma lemma (source0.tex lines 413--422). The article's complete original proof of it is delivered with it (source0.tex lines 423--436, one \texttt{proof} environment of 14 lines). No mathematics is written, repaired or re-derived: the statement and the proof are the article's own lines, in the article's own notation, and no retained line is substituted or re-flowed.  Retained uncounted as the source-local context the proof needs: lines 406--408.  Nothing was omitted from any retained span, so the package carries no omission record.  No retained line contains a citation, so the wrapper carries no bibliography and no bibliography entry.  No retained line contains a figure environment, an image-inclusion command or drawing code, so no image is delivered and the package declares no asset dependency.  Editorial formatting only, in addition: an AMS \texttt{amsart} preamble that restates the article's own declarations and macros used by the retained lines, namely the environments \texttt{prop}, \texttt{lemma} on separate counters, together with the standard packages those lines need.  The retained environments print in this document as Proposition 1, Lemma 1 in order of appearance, whereas the article prints the same items with the numbering of its own full document.  The complete native source of the article as distributed in the arXiv e-print for this exact version is bundled unchanged as \texttt{sources/176673/source0.tex}.
\begin{prop}\label{AllLaplaciansAreSame}
On a K\"ahler Lie Algebroid $(\A, \rho, [\cdot, \cdot], h)$ we have that $\Delta = \Delta_\p + \Delta_{\overline{\p}}$. Moreover, we have that $\Delta_\p = \Delta_{\overline{\p}} = \frac{1}{2}\Delta$.
\end{prop}

\begin{lemma}
    On a Kähler Lie algebroid $(\A, \rho, [\cdot, \cdot],h)$, the following relations hold:
    \begin{enumerate}
        \item $[L,d_\A]=0$,
        \item $[L,\overline{\p}_\A]=0$,
        \item $[L,\Delta_{\overline{\p}}]=0$, 
        \item $[L,\Delta]=0$,
    \end{enumerate}
    where the commutators are regarded in the graded sense.
\end{lemma}

\begin{proof}
    The first relation follows from the fact that $\omega$ is closed, and implies the second relation.  To see the third relation, note that we have
    \begin{align*}
        \Delta_{\overline{\p}}L \alpha &= \overline{\p}_\A\overline{\p}^\dagger_\A L \alpha +  \overline{\p}^\dagger_\A\overline{\p}_\A L\alpha  \\
        &=\overline{\p}_\A(L \overline{\p}_\A^\dagger +i \p_\A)\alpha + \overline{\p}_\A^\dagger L\overline{\p}_\A\alpha \\
        &=L\overline{\p}_\A \overline{\p}_\A^\dagger \alpha + i\overline{\p}_\A\p_\A \alpha + (L \overline{\p}^\dagger_\A+ i\p_\A) \overline{\p}_\A\alpha \\
        &= L\overline{\p}_\A \overline{\p}_\A^\dagger \alpha + i \overline{\p}_\A \p_\A \alpha + L \overline{\p}_\A^\dagger \overline{\p}_\A \alpha + i \p_\A \overline{\p}_\A \alpha \\
        &=L \Delta_{\overline{\p}}\alpha.
    \end{align*}
    The final relation then follows from Proposition \ref{AllLaplaciansAreSame}, as we have
    \begin{equation*}
        [L,\Delta] =2[L,\Delta_{\overline{\p}}]=0.
    \end{equation*}
\end{proof}

\end{document}
